\documentclass{article}
\usepackage[utf8]{inputenc}
\usepackage{hyperref}
\usepackage{fullpage}
\usepackage{amsmath}
\title{subgraph freeness}
\author{talyaa01}
\date{February 2022}

\newcommand{\eps}{\epsilon}
\newcommand{\light}{\textsf{light}}
\newcommand{\heavy}{\textsf{heavy}}
\renewcommand{\ll}{\textsf{$\ell-\ell$}}
\newcommand{\hh}{\textsf{$h-h$}}
\newcommand{\lh}{\textsf{$\ell-h$}}
\newcommand{\davg}{d_{avg}}
\newcommand{\poly}{\textrm{poly}}

\begin{document}

\section{Meeting 20/7/2022}

In this meeting we mainly recalled what we (believe) we have, and may have found on the way something we were missing in our discussions of the lower bound for non-constant $\alpha$ (and one-sided error).

\subsection{Upper bounds}

Let's first recall some notations (which we changed a bit along the way). We let $L$ denote the set of low-degree vertices, whose degree is at most $\alpha/\eps$ (or so), by $H$ the set of high-degree vertices, whose degree is more than $n^{1/2}$ (or so), and by $M$ the set of medium-degree vertices, which are all other vertices.

The definition of $L$ is such that the total number of edges between pairs of vertices that are both in $M\cup H$ is upper bounded by, say, $(\eps/2)m$, and so we can ignore these edges. That is, if the graph is $\eps$-far from being C4-free, then there are $\Omega(\eps m)$ edge-disjoint C4's of the form $M/H$-$L$-$M/H$-$L$.

It looks like we have a proof that when there are $\Omega(\eps m)$ edge-disjoint C4's of the form $H$-$L$-$H$-$L$, then we will catch a C4 by performing $\tilde{O}(n^{1/4}\alpha^{1/2})$ times $\poly(1/\eps)$ random walks of length $2$ (after sampling edges (almost) uniformly to get vertices in $H$ - I don't think we wrote precisely the cost of this part).

For the case that there are $\Omega(\eps m)$ edge-disjoint C4's with more than two low-degree vertices, it looks like we have a simple argument for an upper bound as we aim for (of $\tilde{O}(n^{1/4}\alpha^{1/2})$ times $\poly(1/\eps)$) only when $\alpha \leq n^{1/4}$.  This is based on selecting $\sqrt{d(v)}\poly(1/\eps)$ neighbors from a vertex $v \in L$ (that participates in sufficiently many C4s), and for each selected neighbor that is in $L$,  asking about all its neighbors.

We don't recall having anything for the case that $\alpha$ is bigger (with three $L$s) and there is also the case of almost all C4s being of the form $M$-$L$-$M/H$-$L$.

We also noted that we have $O(n^{1/4})$ for constant $\alpha$, since we can use the $\sqrt{d(v)}\poly(1/\eps)$ bound for C4s that are not with two $H$s

\subsection{Lower Bounds}

First we recalled that we have an $\Omega(n^{1/4})$ two-sided error lower bound for constant $\alpha$. I think it is written somewhere, but let's recall quickly. In the FREE instances we have a bipartite graph with $\Theta(n)$ vertices on one side (the left) and $\Theta(n^{1/2})$ on the other. Every vertex on the left has two unique neighbors on the right. In the FAR instances we pair the vertices on the left, and each pair has the same two neighbors on the right. I guess on the right we need to have less vertices, so that we'll also have a small independent set of ``leftover'' vertices. Can consider other (more elegant?) variants, which we may have though of in the past.

Then we recalled the construction for the case that $\alpha = \Theta(n^{1/2})$, for a one-sided error lower bound of $\Omega(n^{1/2})$ (which is $\Omega(n^{1/4} \alpha^{1/2})$). Here (only have FAR instances), on each side we have $\Theta(n)$ vertices with $\Theta(\sqrt{n})$ perfect matchings between them. As we noted before, this creates parallel edges, so we need to think whether/how to deal with this. The easy part here is to show the lower bound on the number of queries necessary to see a C4  with constant probability, and the more subtle part (which I think maybe we didn't do in full detail in the past) is showing that whp we indeed get a graph that is far from C4-free. I think in the end (of this meeting) we didn't quite get this to work for $\alpha = \Theta(n^{1/2})$ and constant distance, but only for slightly smaller $\alpha$, and since  both parts of the argument are quite similar for general $\alpha$ (or at least, $\alpha \leq n^{1/2}/\log n$ or so), will turn to that.

The construction for general $\alpha$ still has on one side $\Theta(n)$ vertices (let's say $n$ for simplicity), but on the other side it has $n^{1/2}\alpha$ vertices. The vertices on the lhs have degree $\alpha$, and those on the right have degree $n^{1/2}$.
The edges are selected by picking a random matching between entries of an $n\times \alpha$ table (corresponding to the ``neighbor/edge ports'' of the lhs vertices) and the entries of an $n^{1/2}\alpha \times n^{1/2}$ table (corresponding to the ``neighbor/edge ports of the rhs vertices). Here parallel edges will appear even more, but we continue glossing over this.

To show that $\Omega(n^{1/4} \alpha^{1/2})$ queries are necessary to obtain a C4, we actually show that these many queries are necessary for viewing any cycle. Consider any knowledge graph that is obtained after $n^{1/4} \alpha^{1/2}/c$ queries (for some sufficiently large constant $c$). This graph contains at most $2 n^{1/4} \alpha^{1/2}/c$ vertices. So now, if we condition on this knowledge graph and bound the probability that a neighbor query $(v,i)$ will be answered by some vertex $u$ that is in the knowledge graph, we get that it is at most $2 n^{1/4} \alpha^{1/2}/c$ divided by $n^{1/2}\alpha$, which is $2/(cn^{1/4}\alpha^{1/2})$ (we need to be a bit more precise here -- the conditioning on edges already in the graph only reduces the probability of hitting a vertex in the knowledge graph, but there are also non-edges, and these maybe have an opposite effect). For pair queries it looks like it will be something like $\alpha /(n^{1/2}\alpha) = 1/n^{1/2}$.

Let's turn to showing that we get whp a graph that has $\Omega(m)$ disjoint C4s. The first part of the argument is showing that for every edge, the expected number of C4s that it participates in is $\alpha$.
Consider an edge $(u,v)$. For any choice of sets of neighbors $\Gamma(u)$ and $\Gamma(v)$, and for each pair $(u',v') \in \Gamma(v) \times \Gamma(u)$, the probability that we have an edge between them is (roughly) $\alpha \times 1/(n^{1/2}\alpha) = 1/n^{1/2}$. Since we have $n^{1/2} \cdot \alpha$ such pairs, we get  $\alpha$ in expectation. Hence, for any factor of $2^i$, the fraction of edges that participate in more than $2^i\cdot \alpha$ C4s, is at most $1/2^i$, where the maximum number of C4's that an edge can participate in, is $n^{1/2}\cdot \alpha$. We did a simple bucketing argument, with $O(\log n)$ buckets, by which we got that after removing all edges that reside on more than $c' \alpha$ C4s, we removed at most $\log n \cdot n\alpha \cdot \alpha$ C4s, which for $\alpha$ sufficiently smaller than $(n/\log n)^{1/2}$, is much smaller than $n^2$.

Now we use the fact that the expected number of C4s is $n^4 \cdot (1/n^{1/2})^4 = n^2$. I guess here we need to show that whp we won't get significantly less than this expected value (expectation is not enough, and the C4s are not independent of each other, so maybe need Chebishev).

So we should be left with $\Omega(n^2)$ C4s, where each remaining edge resides on at most $c'\alpha$ (maybe $c'$ is simply $ 2$), so if we run a greedy process for $n\alpha/c''$ iterations, where in each iteration we  assign a ``free'' C4 to an edge (and mark all C4s it participates in as ``un-free''), we'll ``use up'' at most $(n \alpha/c'') \cdot c'\alpha$ C4s, which is less than $(c'/c) n^2$.

Other than doing things more carefully and precisely, there is a question as to whether the $\log n$ factor is really necessary.

\section{Meeting 4/8/2022}

We continued in checking out whether applying the ``matching-based construction'' to a ``diluted'' version of the Fox-graph gives us a higher lower bound than we have. By this we mean that we start with a bipartite Fox-graph, and first keep each edge independently with probability $\alpha/\sqrt{n}$ (this of course will create difficulties in a rigorous lower bound argument because of the degrees, but let's ignore this for now, and let's also assume that each vertex has degree $\alpha$ after this). We then add edges (``red'') in the following manner. Each vertex $x$ selects a random matching $\sigma_x$ of its neighbors. Now each $u$ on the lhs of the graph considers each matched pair $v_1,v_2$ of its neighbors (i.e., $\sigma_u(v_1)=v_2$ and $\sigma_u(v_2)=v_1$), and adds the following two edges: $(\sigma_{v_1}(u),v_2)=(\sigma_{v_1}(u),\sigma_u(v_1))$ and
$(\sigma_{v_2}(u),v_1)= (\sigma_{v_2}(u),\sigma_{u}(v_2))$. We thought of these edges as ``being added'' by $u$, but maybe we should view each red edge as being added by an original (``black'') edge. That is, $(u,v)$ adds the edge $(\sigma_v(u),\sigma_u(v))$. We'll see if this comes in useful.

We then started going over all the C4s that we get and a vertex $u$ is part of. These include between 1 and 4 red edges. We went through all those with 1 edge and part of those with 2 edges. It remains to finish going over those with 2 edges and then to check 3 and 4. I won't go over all options (we should find a way to do this in a less case-by-case manner, but not clear yet how). I'll just recall that, on a high level, what we did in the cases that it is not ``a C4 that $u$ is in charge of'' is something along the following lines. I wanted to write this generally, but let's exemplify with the case that there are two red edges in the C4 that are incident to $u$ (see page 6 on the new jamboard).

That is, we have a C4:  $(u,x,y,u')$ where $(u,x), (u,y)$ are red edges and $(x,u'), (y,u')$ are black edges. This means that we have two neighbors of $u$, denoted $v_1,v_2$, and two additional vertices, $u_1$ and $u_2$ such that $u_1$ in a neighbor of $v_1$ and $x$ and $u_2$ is a neighbor of $v_2$ and $y$. Let's consider here the case that all vertices are distinct (this is not necessarily the case, as is also shown on the jamboard). In order to count them, we start with the original Fox graph, We have all $n$ pairs of neighbors $v_1,v_2$ of $u$, times all $n$ vertices $u'$, times all $n$ pairs of neighbors $x,y$ of $u'$, which gives us $n^3$, and determines $u_1$ (the unique common neighbor of $x$ and $v_1$) and $u_2$ (the unique common neighbor of $y$ and $v_2$). Next we have the probability of $(\alpha/sqrt{n})^8 = \alpha^8/n^4$ that all black edges in the octagon remain after the dilution process, and the probability of $\alpha^{-2}\cdot \alpha^{-2} = \alpha^{-4}$ that both red edges are selected (according to the random $\sigma$'s). This gives us $\alpha^4/n$. Since there are $\alpha^4$ pairs of walks of length 2, we will need $n^{1/2}$ random walks to catch such a C4.

We saw that several other C4's also required $n^{1/2}$, and some (like the ``intended'' ones) required $\alpha^{3/2}$. So smearing quite a bit, it may be the case that we can get a lower bound of $\min(n^{1/2},\alpha^{3/2})$ for the random walks algorithm. If this is correct (need to deal with quite a few more details and get rid of ``smearings''), one question is whether we can claim that the random walks algorithm is in general optimal (this is kind of an interesting question in its own right). Another question (even if we remain with random walks), is whether we can have a constructions in which not all vertices have (roughly) the same degree, $\alpha$, that is, not all are in $L$ (low degree), but rather we have C4s with medium-degree vertices. To get this we would need to remove vertices on one side and then also remove edges.

\newpage

\section{Meeting 14/8/2022 (and following it)}
We mainly talked about a special case of the following, which turned out that wouldn't give more than $n^{1/4}\alpha^{1/2}$, so here I am describing the general case, which maybe can give something better with the right setting of parameters. So I am copying something from an email and then adding some thoughts.

Let the set of medium-degree vertices be denoted by $M$. If we want to have on the left side $n$ vertices with degree $\alpha$ (the set $L$), and on the right side have $|M|$ vertices with degree $d$, then $|M| = n\alpha/d$. So, starting with the initial bipartite Fox graph, we independently keep each vertex on the right hand side with probability $\alpha/d$ (recall that $d>\alpha$). At this point vertices in $M$ still have degree $n^{1/2}$ and vertices in $L$ have degree $\alpha n^{1/2}/d$ in expectation. So we next keep each edge independently with probability $d/n^{1/2}$. (Of course all these independent choice give us things only in expectation, and we’ll need to think how to deal with that, but let’s ``flow'' with it.)

Our current lower bound construction (where vertices have degree $\alpha$ (in expectation), will hopefully give us a lower bound of the form $\min(\alpha^{3/2},n^{1/2})$, which is $n^{1/2}$ for $\alpha \geq n^{1/3}$, but lower for $\alpha < n^{1/3}$ (and in particular, not very meaningful for small $\alpha$).
The question is what can we get from the above, and in particular, what should be the setting of $d$. It looks like we'll be Ok in terms of the distance to C4-freeness for every $d$ (didn't check this out, but we saw it is Ok for $d= \alpha$ and for $d=n^{1/2}$). The question is what happens with all the cases we need to analyze for the probability of length-2 random walks colliding.

Before continuing, note that for $\alpha = n^{1/3}$, the bound of $n^{1/2}$ is higher than $n^{1/4}\alpha^{1/2}$, which is $n^{5/12}$ (and in general it is higher for $n^{1/3} < \alpha < n^{1/2}$), and that $\alpha^{3/2}$ is higher than $n^{1/4}\alpha^{1/2}$ for $\alpha > n^{1/4}$, so $\min(\alpha^{3/2},n^{1/2})$ improves on $n^{1/4}\alpha^{1/2}$ for $n^{1/4} < \alpha < n^{1/2}$.

Returning to the attempt at a lower-bound construction with vertices having ``medium-degree'' $d$, let's look at the cases with a single red edge. Among them, let's first look at the cases of catching a C4 $(u,v,u',v')$ where $(u',v')$ is the red edge, and we start from either $u$ or $v$, where $u \in L$ and $v\in M$.
\begin{enumerate}
\item One case is when $(u',v') = (\sigma_v(u), \sigma_u(v))$. To catch such a C4 from $u$, we would need $((\alpha d)^2/\alpha)^{1/2} = d\alpha^{1/2}$ length-2 random walks, while from $v$ we'd need $((\alpha d)^2/d)^{1/2} = d^{1/2}\alpha$, which is smaller (since $d \geq \alpha$).
\item The other case is when $(u',v') = (\sigma_{v_1}(u_1), \sigma_{u_1}(v_1))$ for another  edge $(u_1,v_1)$ (where $u_1$ is a neighbor of $v'$ and $v_1$ is a neighbor of $u'$). That is, we have a ``hexagon'' with a red-edge in the midle. The number of hexagons that $u$ belongs to (similarly, $v$) in the original Fox graph is $n^2$. But now begins the difference between $u$ and $v$.

    For $u$, the probability that such a hexagon ``survives the diluting'', is $(\alpha/d)^3 \cdot (d/\sqrt{n})^6$ (the first term is due to $v, v', v_1$ ``surviving'' and the second is due to all edges surviving), while for $v$ it is $(\alpha/d)^2 \cdot (d/\sqrt{n})^6$. For both we then multiply by $1/(\alpha d)$ (the probability of the red edge). This amounts to $(\alpha^2 d^2)/n$ for $u$ and $(\alpha d^3)/n$ for $v$, which is larger by a factor of $d/\alpha$. Therefore, for $v$ we shall need $\sqrt{n \alpha/d}$.
\end{enumerate}
Comparing the two expressions, we ask when are they equal (and condition on $d\geq \alpha$).
We get that
$d^{1/2}\alpha = \sqrt{n \alpha/d}$, when $d=\sqrt{n/\alpha}$. This is at least $\alpha$ for $\alpha \leq n^{1/3}$.
Plugging in this setting of $d$, we get $n^{1/4}\alpha^{3/4}$. Note that for $\alpha \leq n^{1/3}$ this is upper bounded by $n^{1/2}$, and is at least $\alpha^{3/2}$.

So maybe we can get $\min(n^{1/4}\cdot\alpha^{3/4},n^{1/2})$. We should of course check all cases for this. I think that for the other cases of a single red edge (incident to $u$ or incident to $v$) we get the same expression as in the second item above. In all cases it looks like starting from $v$ will always give a smaller lower bound, so we can focus only on that. I also checked the case in which there are two edges, which are not incident to $v$ (i.e., in $(u,v,u',v')$, the edges $(u,v')$ and $(v,u')$ are red), and these edges are ``due to'' $(u_1,v_1)$ and $(u_2,v_2)$, where all vertices are distinct (so that we have an octagon in the original Fox graph). I got $\sqrt{n\alpha/d}$ here too.

A few thoughts / simple observations regarding the expression $\min(n^{1/4}\cdot\alpha^{3/4},n^{1/2})$. First, maybe as already noted above, it is $n^{1/4}\cdot\alpha^{3/4}$ for $\alpha \leq n^{1/3}$ and it is $n^{1/2}$ for $\alpha \geq n^{1/3}$. Focusing on the first term, if a vertex $v$ has degree at most $\sqrt{n/\alpha}$ (and it participates in relatively many C4s) then we can catch it with $\sqrt{d(v)}\alpha \leq n^{1/4}\alpha^{3/4}$ length-2 walks. Also note that the set of vertices with degree greater than $\sqrt{n/\alpha}$ has size at most $n\alpha/\sqrt{n/\alpha} = n^{1/2}\alpha^{3/2}$, so that $n^{1/4}\alpha^{3/4}$ is the square-root of this expression. Can we get complexity that is the square-root of the number of vertices on the rhs?

\newpage

\section{Meeting 15/9/2022}

We first went over some more details of the lower bound (for length-$2$ walks). In particular, we discussed one type of C4 $(v,u,v',u')$ with two red edges ($(v,u')$ and $(u,v')$) but in which the black (original Fox) edges ``responsible'' for these edges, $(u_1,v_1)$ and $(u_2,v_2)$ are not both over disjoint vertices. Rather, $v$ equals one of $v_1$ or $v_2$, say, $v_2$, so that the underlying black structure is a C6 plus one edge. We said that if we fix $v$, then in the original Fox graph we have at most $n^2 \cdot n^{1/2}$ options for the black-edges structure ($n^2$ takes into account the choices of $v'$ and $v_1$, which together with $v$ determine $u_1$ and $u'$, and $n^{1/2}$ is the number of choices of $u$). Now we multiply by $(\alpha/d)^2$ (the probability that both $v'$ and $v_1$ survive the diluting of vertices), and by $(d/n^{1/2})^7$ (the probability that all seven black edges survive the diluting of edges). Finally, we multiply by $(1/(d\alpha))^2$ (the probability that $(v_1,u_1)$ ``added'' $(v,u')$ and $(v,u_2)$ added $(u,v')$). This gives us $d^3/n$. So we get a lower bound of $\left(\frac{(d\alpha)^2}{d^3/n}\right)^{1/2} = n^{1/2}\alpha/d^{1/2}$ which is bigger than $n^{1/4}\alpha^{3/4}$ for $d = (n/\alpha)^{1/2}$.  We also recalled that for the case in which $v_1$ and $v_2$ are different from $v$ (and I guess we also need to check $u_1$ and $u_2$ being different from $u$), we got
$n^{1/2}\alpha^{1/2}/d^{1/2}$ .


We recalled which cases are easy for us, so that we can focus on the non-easy one (in the attempt to match the lower bound of $\Omega(\min(n^{1/4}\alpha^{3/4},n^{1/2}))$). Actually, this is written in the previous paragraph, for $\alpha \leq n^{1/3}$, and for $\alpha > n^{1/3}$ the easy case is when $d$ is not much larger than $\alpha$. So we need to focus on the other cases.

We then talked about there being ``many'' C6s when either $\alpha \leq n^{1/3}$ and $d > \sqrt{n/\alpha}$ so that $d^2 \alpha > n$, or $\alpha > n^{1/3}$ (and just $d> \alpha$), so that also $d^2\alpha > n$. We tried to understand whether combining this with there being many disjoint C4s means there must be many additional C4s, but I think we didn't get very far with this - I guess we saw that it is possible that we won't get unintentional C4s based on one starting vertex $v$.

\bigskip
The last thing we talked about and want to think about more is the following.
Given a bipartite graph $G$, whose left side is denoted by $L$, with vertices having degree $\alpha$ and whose right side is denoted by $M$, and whose vertices have degree $d > \sqrt{n/\alpha}$ for $\alpha \leq n^{1/3}$, or just $d> \alpha$ when $\alpha > n^{1/3}$, we considered defining am $|M| \times |M|$  ``wedges'' matrix as follows. In each entry $(v_i,v_j)$ of the matrix (corresponding to a pair of vertices in $M$), we put all vertices $u_k$ of $L$ such that $(v_i,u_k,v_j)$ is a wedge.

What do we know about this matrix, and what do we hope we can get from it?
First, for each vertex $u_k \in L$, the entries of the matrix that it belongs to correspond to a ``square sub-matrix'' of size $\alpha \times \alpha$ (whose rows and columns correspond to the neighbors of $u_k$ in $M$). This is the constraint imposed by the matrix being defined by a graph. We also know  that for each row $v_i$ (similarly columns), the sum of the sizes of its entries is $d\alpha$, and there are $d$ distinct vertices $u_k$ that appear in the entries in this row. Finally (in terms of what we know), suppose that for each row we do the following. We consider all entries of size at least two. Now, we color each $v_k$ that appears in one of these entries by red at most once, if it has a partner $v_{k'}$ in the same entry that we can match it to (this is not very formal, but hopefully sufficiently clear). Then we know that there are $\Omega(\eps m)$  red occurrences of vertices from $L$ in the matrix, and actually, that in each row we have $\Omega(\eps' d)$ such red occurrences.

What we would like to lower bound is the sum, over all entries of size at least 2, of the size of the entry squared.
(In retrospect, this is not precise, but we would like there to be relatively many rows for which this sum is relatively large).

\section{Meeting 21/9/2022}

First I am enclosing below the email that I sent before the meeting:

In our last meeting we weren't sure whether the ``red markings” in the matrix we discussed should be on edges or wedges and exactly how this should be represented in the matrix. Maybe it will be helpful to think about the matrix determined by the subgraph that contains only the $\Omega(\eps m)$ edge-disjoint C4s. In the actual graph there may be additional (“black”) edges, but a special case is on in which there are no extra edges ($\eps$ is a constant), and also, hopefully the extra edges do not make our life (much) harder.

In fact, let’s first fill only those entries in the matrix that correspond to the C4s. That is, if we denote the matrix by $W$, and we have a C4 $(v_1,u_1,v_2,u_2)$, then $W[1,2] = \{u_1,u_2\}$. So initially, in each row $i$, there are entries of even size, with distinct elements, and the sum of the sizes of these entries is $d$ (or $\eps' d$). (I initially wrote that each entry is of size two, but this is not necessary correct: it is possible that we have the following two edge-disjoint C4s: $(v_1,u_1,v_2,u_2)$ and $(v_1,u_3,v_2,u_4)$, so that $W[1,2] = \{u_1,u_2,u_3,u_4\}$.)
% there are $d$ (or maybe $\eps’ d$) columns $j$ such that $W[i,j]=\{u_k,u_\ell\}$, and the pairs in these entries are all distinct.

What we would like to show is that due to the ``(generalized) square sub-matrices constraint” (i.e., that if $u$ belongs to an entry $W[i,j]$ and to an entry $W[i',j’]$, then it also belongs to $W[i,j’]$ and $W[i',j]$), the completely filled matrix as we defined it, must contain many C4s (this ``many” is not well defined, but I think can be formalized, and is related, as we said, to the sum over all entries $W[i,j]$ of $W[i,j](W[i,j]-1)$). Possibly we’ll also want to use the fact that each  vertex $u \in L$ participates in relatively many disjoint C4s. This means that after the first step of filling the matrix, each $u\in L$ belongs to roughly $\alpha$ entries $W[i,j]$, where no two of them are in the same row/column (in addition to what is written above about each row).

For example, suppose we have one C4, $(v_1,u_1,v_2,u_2)$ and another (edge-disjoint), $(v_3,u_1,v_4,u_2)$, then initially, in the matrix, we write $\{u_1,u_2\}$ in entries $[1,2]$ and $[3,4]$ (it might be helpful to also fill in $[2,1]$ and $[4,3]$ in order to better visualize the (generalized) square sub-matrices. But then, $u_1$ and $u_2$ must both also belong to the entries $[1,3],[1,4],[2,3],[2,4]$ (as well as $[3,1],[4,1],[3,2],[4,2]$). This is an example where two disjoint C4s induce additional C4s, and we also can have three C4s or four C4s inducing additional C4s.

\bigskip
We then turned to discuss what the $W$ matrix looks like in the Fox graph (possibly diluted version). We also started filling the matrix by entries that correspond to wedges (centered at  vertices from $L$) where each  vertex in $L$ has at most one entry in each row and column and in each row there are $d$ entries, each of size 1. So while in the initial step of filling the $W$ matrix in the case of having many disjoint C4s we have even-sided entries (of size at least 2), here the entries are of size 1. Furthermore, we know that if we now continue filling $W$ starting from this initial stage, then we'll get addition entries of size 1, but no entry of larger size.

This led us to think about getting a different lower bound construction. We first observed that if we just modify each entry $W[i,j] = \{k\}$ to $W[i,j] = \{k,k'\}$ (where this is the same $k'$ in each entry and $k$ belongs to), then we get a lower lower bound than we have. This construction corresponds to taking a Fox graph (once again, possibly after diluting), and adding a ``twin'' vertex $k'$ for each  vertex $k$ in $L$, where $k'$ is a twin in the sense that it has exactly the same  neighbors in $M$ as $k$. Consider any  vertex $v \in M$. Among the $(\alpha d)^2$ pairs of length-2 paths starting from $v$, there are $\alpha d$ pairs that close a C4 (for each path $(v,u,v')$, where there are $\alpha d$ such paths, there is exactly one path $(v,u',v')$, that closes a C4 with the first, where $u'$ is the twin of $u$). So we shall need $\Omega\left(\left(\frac{(\alpha d)^2}{\alpha d}\right)^{1/2}\right) = d^{1/2}\alpha^{1/2}$ which is upper bounded by $n^{1/4}\alpha^{1/2}$.

So we thought about how to modify the above. In particular, we considered taking two Fox graphs (before diluting) that share the same set of right-hand-size vertices (which after diluting will be the set $M$) but have different sets of low-degree ($L$) vertices. We said the second graph will be defined by taking a permutation over the names of the rhs vertices, but it looks like we didn't use this. We then perform the diluting. Assuming we don't really need the random permutation, we might as well start with each  vertex $u \in L$ having a twin ($u' \in L'$). Note that this is different from what we did above, since we are first adding a twin and then diluting, while before we first diluted and then added a twin. Also note that in the special case where there is no diluting, so that $\alpha = n^{1/2}$, everyone has degree $n^{1/2}$ and the constructions are the same, we get a lower bound of $\Omega(n^{1/2})$, which is fine.

So now let's see what happens when we dilute. We discussed together the case that $\alpha \leq n^{1/3}$ where we dilute both vertices and edges. Reut and I later verified the case that $\alpha > n^{1/3}$, where we only dilute edges. Let's actually start from that, since it is a bit simpler.

Before diluting, every pair of rhs vertices have exactly two common neighbors - one from each of the initial Fox graphs. There are no other C4s. Hence, we initially have $n^2$ C4s. The probability that each such C4 survives is $(\alpha/n^{1/2})^4$, and we get an expected number of $\alpha^4$ C4s (recall that $\alpha \geq n^{1/3}$, so that this is at least $n\alpha$). On each edge, the expected number of surviving C4s is $n^{1/2} \cdot (\alpha/n^{1/2})^3 = \alpha^3/n$, so that we expect to have $\alpha^4 / (\alpha^3/n) = n\alpha$ edge-disjoint C4s. What about the number of length-2 paths that are necessary for catching a C4? If all is as in expectation, we get a lower bound of $((\alpha \cdot \alpha)^2 / (\alpha^4 /n)^{1/2}$ (recall that here $d=\alpha$), which is $n^{1/2}$.

In the case that $\alpha \leq n^{1/3}$, we also initially have $n^2$ C4s. Here the probability that such a C4 survives is $(\alpha /d)^2 \cdot (d/n^{1/2})^4$, we get an expected number of $\alpha^2 d^2$ C4s, which for $d \geq (n/\alpha)^{1/2}$ is at least $n\alpha$. On each edge, the expected number of surviving C4s is $n^{1/2} \cdot (\alpha/d)\cdot  (d/n^{1/2})^3 = d^2\alpha/n$, so that we expect to have $\alpha^2 d^2 / (d^2\alpha/n) = n\alpha$ edge-disjoint C4s. What about the number of length-2 paths that are necessary for catching a C4? Since we have (roughly) $n\alpha/d$ vertices in $M$, we have an average of $(\alpha^2 d^2)/(n\alpha /d) = d^3\alpha/n$ C4s for each $v\in M$. If all is as in expectation, we get
$((\alpha d)^2 / (d^3\alpha/n )^{1/2} = (n \alpha /d)^{1/2}$ (which is $|M|^{1/2}$), and for $d = (n/\alpha)^{1/2}$ (which is the smallest $d$ we can take so as to get a graph that is far from C4 free), we get $n^{1/4}\alpha^{3/4}$.

So it seems like this gives us the same lower bound that we had, but in a simpler way (without lots of cases).


At this point we can either try to increase this lower bound by a different construction, or return to the question of the upper bound, by analyzing the $W$ matrix based on any ``initial filling'' that satisfies the constraints we wrote above.

\iffalse
\section{Following our meeting on 29/9/2022 and the email correspondence}

Let $W$ be the matrix defined above after it is initially filled based on the set of disjoint C4s. We said that it looks like we can assume that the pairs in the matrix can be ordered so that in each one is ``first'' and the other ``second'', where the subset of first vertices is disjoint from the second. We should double check the details, but hopefully the argument should be something like the following. Starting from the actual initial $W$, for each vertex from $L$ that participates in the matrix, we flip a coin to decide if it wants to be first or second. We then keep all pairs that made a different choice. The expected number of such pairs is half the initial number, and furthermore, for each vertex, in expectation, it will have half of its partners. We need to verify that this is enough - in terms of the total number, it is, since there exists an assignment that gives at least the expected value, but individually, it is not so (maybe not so important?)
\fi

\section{Trying to define a greedy algorithm for disjoint C4s (after meeting on 12/10)}

The following starts like what we talked about, but then continues a bit differently.

On a high level, in each step
we have a bipartite subgraph in which the degree of every vertex is even, the edges can be partitioned into disjoint C4s, and there are no other C4s. A subset of the vertices are ``closed'', in the sense that we set all their incident edges. In each step we close one more vertex from the right hand side. There is  an order on the vertices from each side, where the order is just according to the step in which a new vertex is added.

Given a vertex that we want to close (where we take the next vertex that is not yet closed on the right hand side, according to the above order), we do the following. We add its incident edges one by one, where the rule is that for each edge, we pick the first vertex on the other side such that adding this edge will not close a C4 with any existing edge that is already assigned to a C4. If needed, we add new vertices. We match the endpoints of these edges according to the order the edges were added. For each matched pair, we add two edges to a common endpoint such that here too we don't get a C4 with an edge that is already assigned to a C4.

In particular:
\begin{enumerate}
\item In the first step, we start with vertex $R_1$ (first vertex on the right). We set its neighbors to be $L_1,\dots,L_d$ (I am writing $d$ probably as a shorthand for $\sqrt{n}$). We now match them and for each matched pair, we put edges to a new vertex on the right. Let's call these vertices $R_2,\dots,R_{d/2}$, where $R_2$ is the common neighbor of $L_1$ and $L_2$, $R_3$ is the common neighbor of $L_3$ and $L_4$, and so on (we may consider naming them differently).
\item We next close $R_2$. Note that it cannot neighbor any vertex among $L_3,\dots,L_d$, since if it did, the edge $(R_1,L_1)$, which is already assigned to a C4, would participate in the C4 $(R_1,L_1,R_2,L_i)$. So we give $R_2$ $d-2$ neighbors that are new vertices on the left hand side. We match them, and for the pairs we get, we can actually use $R_3,\dots,R_{d/2}$ as common neighbors.
\item Continuing with $R_3$ (who at this point has degree 4, so it needs $d-4$ additional neighbor, it can't share a neighbor with $R_1$ by the same reasoning that $R_2$ couldn't. It also can't share neighbors with $R_2$, because it already shares two neighbors with $R_2$, so it must have new neighbors. Now, if I am not mistaken, after pairing these $d-4$ neighbors, we can close C4s with the same vertices that $R_1$ and $R_2$ did (e.g., the third and fourth neighbors of $L_3$ (which are its first two new neighbors) will be neighbors of $L_4$).
\end{enumerate}
If we continue in this manner, it looks like in order to close the remaining $R_i$s that we currently have, we need to add for each $R_i$, $d-2i$ new neighbors on the left (when we read it, it has $2i$ neighbors), and we can close C4s with $R_{i+1},\dots,R_{d/2}$, so that we don't add any new neighbors on the right.

Maybe we should call the above, the first phase. At the end of this phase we have $d/2$ closed vertices on the right, and $\Theta(n)$ (for $d=\sqrt{n}$) vertices on the left, where each vertex on the left has exactly two neighbors. The vertices on the left come in pairs, where each pair has their distinct two common neighbors on the right (distinct, but not disjoint). In the second phase we add a new vertex on the right. When we select its neighbors, it is not allowed to take any two vertices that have a common neighbor on the right. If we think of each set of edges added at a stage of the first phase as a \emph{batch}, then it can take at most one vertex from each batch. As we shall see, it is actually more restricted. We then add vertices on the left so that this new vertex will have degree $d$.
% This will mean that it has degree $d/2$, and we will also have to take into account that the batches have decreasing sizes. We can either add vertices on the left, or see what happens if we don't, but just allow vertices to have degree $\Theta(d)$ instead of $d$ (and maybe even smaller degree on the left than on the right).

Trying this, in the second phase,  let's call the new vertex we add on the right $R_{2,1}$ (and maybe rename the $R_i$s from above to be $R_{1,i}$). If it takes $L_1$ as its first neighbor, then, since the neighbors of $L_1$ are $R_1 = R_{1,1}$ and $R_2 = R_{1,2}$, it can take the first neighbor of $R_3$ that is not already a neighbor of $R_{1.1}$ or $R_{1,2}$. In general, let's denote by $\Gamma^{*}(R_{1,i})$ the neighbors of $R_{1,i}$ that are not neighbors of $R_{1,i'}$ for $i' < i$. It then can take the first neighbor in $\Gamma^*(R_5)$, $\Gamma^*{R_7}$ and so on. This gives it $d/4$ (or so) neighbors, so we give it $3d/4$ additional new neighbors. I am not sure if it important how we match them, so let's say for now, according to this order, and then we close C4s for every matched pair, with a new vertex on the right, where we name these new vertices, $R_{2,2},\dots,R_{2,d/2}$.
At the end of this phase, these vertices are closed. As for the vertices on the left, the vertices in $\Gamma^*(R_{1,i})$ for odd $i$ have degree 4, and all other vertices have degree 2. This raises the question if we shouldn't have used vertices in $\Gamma^*(R_{1,i})$ for even $i$. Maybe we can/should do this in the next phase, but then we need to define the greediness in a consistent manner. One thing we could take into account when choosing neighbors is the restrictions they induce (which is roughly linear in the number of neighbors they already have). Such a rule should cause the algorithm to select neighbors from $\Gamma^*(R_i)$ for even $i$ in the next phase.
% Say it selects the first vertex in each batch to be its neighbor, and it pairs them according to this order. Since the other vertices currently on the right are already closed, for each pair it adds a new vertex $R_{2,i}$, for $i= 2,\dots,d/4$. Actually, this doesn't seem to be heading in the right direction if in each phase $j$ we add $d/2^j$ new vertices on the right.
%
% So perhaps we should add for $R_{2,1}$, $d/2$ additional neighbors (``with no strings attached''). Pairing both ``the old neighbors'' and the ``new ones'' (separately?) we then close $d/2$ C4s with new vertices on the right, similarly to what we did in the first step of the first phase.

%The question is whether we can formulate some invariants that we have, which ensure that this process ends successfully.
% Alternatively, we can start by dealing with $R_1$ and $L_1$ together, so that $L_1$ also has edges to $R_2,\dots,R_d$, we have an edge between $L_2$ and $R_2$, so that $(R_1,L_1,R_2,L_2)$ is a C4, and both on the right hand side and the left hand side, we have additional $(d-2)/2$ vertices, that close C4s with either $R_1$ and its paired neighbors, or $L_1$ and its paired neighbors.
%
% In both cases, the question is what is the next vertex to close. In the second case it seems natural to deal with $L_2$ and $R_2$. But what about the first case? And also, in the second case, with what vertex do we continue. Does it make sense, in the first case, to continue by closing vertices on the rhs, and when finished with them return to the lhs (or actually, if all vertices on the rhs are closed, then we should be done)?

\section{Sketches for what we have on C6}

Very shortly, it looks like we have a lower bound of $\Omega(n^{1/3})$ for constant $\alpha$, which follows from the lower bound of $\Omega(\min\{d,n/d\})$ for testing triangle-freeness, i.e., C3-freeness (where $d$ is the average degree), when applied to $d=n^{1/2}$. The lower-bound follows by a reduction from C3 to C6, where each edge is replaced by a path of length-2. (A similar construction can be applied to reduce to $C_{3t}$ for any integer constant $t$.) In addition to showing that distance is maintained in this reduction, the important thing to note is that we can efficiently answer queries to the resulting graph (which has $n'=n+m$ vertices and $m'=2m$ edges), given access to the original graph (in the triangle-free lower bound construction), because the original graph is regular.

CONTINUED IN SEPARATE FILE


\end{document}
